% ============================================================================
% PLANTILLA 05 — PRÁCTICA DIRIGIDA
% ----------------------------------------------------------------------------
% Para: sesiones guiadas — ejercicios modelo resueltos paso a paso
% (entorno {desarrollo}, siempre visible) seguidos de ejercicios propuestos
% para el estudiante.
% Compilar: latexmk -pdf plantilla-05-practica-dirigida.tex
% ============================================================================
\documentclass{academic-exam}

\universidad{Universidad Nacional de San Cristóbal de Huamanga}
\facultad{Facultad de Ciencias Económicas, Administrativas y Contables}
\escuela{Escuela Profesional de Economía}
\curso{Econometría I}
\codigocurso{EC-352}
\docente{Mg. Nombre Apellido}
\periodo{Semestre 2026-I}
\tipoevaluacion{Práctica Dirigida 03}
\fechaevaluacion{21 de julio de 2026}
\duracion{2 horas}
\modalidad{Presencial · grupal}

\begin{document}
\membrete

\begin{instrucciones}[Objetivo de la sesión]
Al finalizar la práctica, el estudiante estima e interpreta un modelo de
regresión lineal simple por mínimos cuadrados ordinarios (MCO), evalúa la
bondad de ajuste y realiza inferencia sobre los coeficientes.
\end{instrucciones}

\begin{definicion}[Estimador MCO]
En el modelo $y_i = \beta_0 + \beta_1 x_i + u_i$, los estimadores de
mínimos cuadrados ordinarios son
\[
  \hat{\beta}_1 = \frac{\sum_{i}(x_i-\bar{x})(y_i-\bar{y})}
                       {\sum_{i}(x_i-\bar{x})^{2}},
  \qquad
  \hat{\beta}_0 = \bar{y} - \hat{\beta}_1\bar{x}.
\]
\end{definicion}

% ==================================================== EJERCICIO MODELO =====
\section*{\color{colorprimario}A · Ejercicio modelo}

\begin{ejercicio}[][Regresión simple paso a paso]
Con una muestra de 8 departamentos se tiene el gasto en educación por
alumno ($x$, miles de soles) y el rendimiento promedio en la ECE ($y$,
puntos estandarizados):

\begin{datos}
\centering
\begin{tabular}{@{}l *{8}{S[table-format=3.1]}@{}}
  \toprule
  $x_i$ & 2.1 & 2.8 & 3.0 & 3.4 & 3.9 & 4.2 & 4.8 & 5.0\\
  $y_i$ & 480.0 & 495.0 & 505.0 & 512.0 & 520.0 & 528.0 & 541.0 & 547.0\\
  \bottomrule
\end{tabular}
\end{datos}

Estime por MCO la recta $\hat{y} = \hat{\beta}_0 + \hat{\beta}_1 x$ e
interprete los coeficientes.

\begin{desarrollo}
\textbf{Paso 1 — Medias.} $\bar{x} = 29{,}2/8 = 3{,}65$;\quad
$\bar{y} = 4\,128/8 = 516$.

\textbf{Paso 2 — Sumas de cuadrados y productos cruzados.}
\[
  S_{xx} = \sum (x_i-\bar{x})^{2} = 7{,}46, \qquad
  S_{xy} = \sum (x_i-\bar{x})(y_i-\bar{y}) = 166{,}9.
\]

\textbf{Paso 3 — Estimadores.}
\[
  \hat{\beta}_1 = \frac{S_{xy}}{S_{xx}} = \frac{166{,}9}{7{,}46}
  = 22{,}37,
  \qquad
  \hat{\beta}_0 = 516 - 22{,}37\,(3{,}65) = 434{,}33.
\]

\textbf{Paso 4 — Interpretación.} Un incremento de mil soles en el gasto
por alumno se asocia con un aumento promedio de 22,4 puntos en el
rendimiento; el intercepto (434,3) es el rendimiento esperado con gasto
nulo, fuera del rango muestral y sin lectura económica directa.
\end{desarrollo}
\end{ejercicio}

\begin{observacion}
La pendiente MCO solo admite lectura \emph{causal} bajo el supuesto de
exogeneidad, $E(u\mid x)=0$. Con datos observacionales como estos, la
correlación puede reflejar factores omitidos (ruralidad, pobreza,
gestión educativa).
\end{observacion}

% ================================================ EJERCICIOS PROPUESTOS ====
\section*{\color{colorprimario}B · Ejercicios propuestos}

\begin{ejercicio}[][Bondad de ajuste]
Con los datos del ejercicio modelo, calcule la suma de cuadrados
explicada, el coeficiente de determinación $R^{2}$ y el error estándar de
la regresión. Interprete cada resultado.
\lineasrespuesta{7}
\end{ejercicio}

\begin{ejercicio}[][Inferencia]
Contraste $H_0\colon \beta_1 = 0$ contra $H_1\colon \beta_1 \neq 0$ al
5\,\% de significancia. Reporte el estadístico $t$, el valor crítico y su
conclusión.
\lineasrespuesta{7}
\end{ejercicio}

\begin{ejercicio}[][Predicción]
Prediga el rendimiento esperado para un gasto de 4,5 mil soles y comente
qué precauciones exige extrapolar fuera del rango muestral
$[2{,}1;\,5{,}0]$.
\lineasrespuesta{6}
\end{ejercicio}

\begin{observacion}[Para la siguiente sesión]
Replique la estimación en R con \texttt{lm(y \textasciitilde{} x)} y
traiga impresa la salida de \texttt{summary()}. Compararemos los cálculos
manuales con la salida computacional.
\end{observacion}

\findelexamen
\end{document}
